% 妊神星（综述）
% license CCBYSA3
% type Wiki

本文根据 CC-BY-SA 协议转载翻译自维基百科\href{https://en.wikipedia.org/wiki/Haumea}{相关文章}。

Haumea（小行星编号：136108 Haumea）是一颗位于海王星轨道之外的矮行星。$^{[25]}$
它由加州理工学院（Caltech）的 Mike Brown 领导的团队于 2004 年在帕洛马天文台发现；随后，西班牙内华达山脉天文台 José Luis Ortiz Moreno 领导的团队在 2005 年正式宣布了这一发现，他们在 2003 年拍摄的先前图像中识别出了该天体。从那次宣布起，它获得了临时编号 2003 EL61。

2008 年 9 月 17 日，它被正式命名为 Haumea，源自夏威夷生育与分娩女神。国际天文学联合会（IAU）基于其可能具备矮行星地位的预期而赋予这一名称。根据标称估计，它是仅次于厄里斯和冥王星的第三大已知海王星外天体，其大小约等同于天王星的卫星泰坦妮娅。Haumea 的最早“预发现”（precovery）影像可以追溯至 1955 年 3 月 22 日。$^{[9]}$

Haumea 的质量约为冥王星的三分之一、地球的 1/1400。虽然尚未直接观测到其形状，但根据其光变曲线的计算结果，与 雅可比椭球体（Jacobi ellipsoid）一致（这是一种在满足矮行星条件时会呈现的形状），其长轴大约是短轴的两倍。

2017 年 10 月，天文学家宣布在 Haumea 周围发现了一套环系统——这是首个在海王星外天体以及矮行星上发现的环系统。

Haumea 的引力直到最近仍被认为足以使其达到流体静力平衡，但目前这一点仍不确定。Haumea 修长的形状、快速自转、环结构以及高反照率（源自其晶质水冰表面），被认为都是一次巨大碰撞的结果。这场碰撞使 Haumea 成为一个碰撞族群（Haumea 家族）的最大成员，该族群包括多个大型海王星外天体，以及 Haumea 的两个已知卫星：Hiʻiaka 与 Namaka。
\subsection{历史}
\subsubsection{发现}
由加州理工学院（Caltech）的 Mike Brown、耶鲁大学的 David Rabinowitz，以及位于夏威夷双子天文台的 Chad Trujillo 组成的团队，于 2004 年 12 月 28 日在他们拍摄的2004 年 5 月 6 日的图像中发现了 Haumea。2005 年 7 月 20 日，他们在网上发布了一篇摘要，准备在 2005 年 9 月的一次会议上正式宣布这一发现。$^{[26]}$

在此期间，西班牙内华达山脉天文台安达卢西亚天体物理研究所（Instituto de Astrofísica de Andalucía）的José Luis Ortiz Moreno及其团队，在2003 年 3 月 7–10 日拍摄的图像中找到了 Haumea。2005 年 7 月 27 日晚，Ortiz 向小行星中心（MPC）发送了关于他们发现的邮件。$^{[27]}$

Brown 最初承认发现权应归 Ortiz 团队，$^{[28]}$ 但在得知西班牙团队在官方发现宣布的前一天访问过 Brown 的观测日志后，开始怀疑对方存在欺诈行为，而 Ortiz 团队在宣布中并未像惯例那样披露这一事实。这些日志包含足够的数据，使 Ortiz 团队能够在其 2003 年的图像中提前定位（precover）Haumea，并且在 Ortiz 在7 月 29 日安排望远镜时间以获取确认图像、并向 MPC 再次提交发现报告之前，又再次访问了这些日志。Ortiz 后来承认访问过 Caltech 的观测日志，但否认任何不当行为，声称自己只是核实他们是否发现了一个新的天体。$^{[29]}$

根据 IAU（国际天文学联合会）规定，小行星的发现权归属于最先向 MPC（小行星中心）提交足够轨道数据以确定其轨道的团队，并且获得发现权的团队有权优先命名该天体。
然而，IAU 在2008 年 9 月 17 日发布的关于 Haumea 命名的官方公告（该命名由专门负责预期为矮行星的天体设立的联合委员会决定）中，并未提及发现者。公告中将发现地点列为西班牙团队的内华达山脉天文台，$^{[30][31]}$ 但最终采用的名称“Haumea”却是 Caltech 团队的提议。

Ortiz 团队曾提议命名为“Ataecina”，取自古伊比利亚春之女神；$^{[27]}$ 但由于该女神属于冥界神（chthonic deity），这一名字更适合用于命名“冥卫星类（plutino）”天体，而 Haumea 并非此类天体，因此未被采用。
